\chapter{Mengen von Mengen}

\section{Mengige Familien und ihre Vereinigung}

\begin{definition}[Mengenfamilie]
Eine \emph{Mengenfamilie} $\mathcal F$ ist hier einfach eine Menge,
deren Elemente als die betrachteten Mengen aufgefaßt werden.
Die Elemente von $\mathcal F$ heißen ihre \emph{Glieder}.
Eine Indizierung gehört nicht zur Definition.
Eine \emph{Teilfamilie} ist eine Teilmenge von $\mathcal F$.
\end{definition}

\begin{definition}[Vereinigung einer Familie]
Die durch das Vereinigungsaxiom bestimmte Menge $\bigcup\mathcal F$
erfüllt
\[
 x\in\bigcup\mathcal F
 \Longleftrightarrow \exists A\in\mathcal F\,(x\in A).
\]
Diese Definition gilt auch für $\mathcal F=\varnothing$.
\end{definition}

\begin{theorem}[Vereinigung, Enthaltensein und leere Familie]
\[
\begin{gathered}
 \bigcup\varnothing=\varnothing,\qquad
 \bigcup\{A\}=A,\qquad \bigcup\{A,B\}=A\cup B,\\
 \bigcup\mathcal F\subseteq U
   \Longleftrightarrow \forall A\in\mathcal F\,(A\subseteq U),\\
 \bigcup\mathcal F=\varnothing
   \Longleftrightarrow \forall A\in\mathcal F\,(A=\varnothing).
\end{gathered}
\]
Für jedes Glied $A\in\mathcal F$ gilt $A\subseteq\bigcup\mathcal F$.
Jede Menge $A$ läßt sich außerdem als Vereinigung ihrer
Einzelmengen schreiben:
\[
 A=\bigcup\{\{a\}\mid a\in A\}.
\]
\end{theorem}

\begin{theorem}[Teilfamilien und Zusammenfassen von Familien]
\[
\begin{gathered}
 \mathcal F\subseteq\mathcal G
      \Rightarrow\bigcup\mathcal F\subseteq\bigcup\mathcal G,\\
 \bigcup(\mathcal F\cup\mathcal G)
       =\bigcup\mathcal F\cup\bigcup\mathcal G,\\
 \bigcup(\mathcal F\cup\{A\})=\bigcup\mathcal F\cup A.
\end{gathered}
\]
Für jede Menge $\mathfrak F$ von Mengenfamilien gilt
\[
 \bigcup\bigl(\bigcup\mathfrak F\bigr)
       =\bigcup\{\bigcup\mathcal F\mid\mathcal F\in\mathfrak F\}.
\]
\end{theorem}

\begin{theorem}[Vereinigung bei Schnitt und Differenz von Familien]
Für beliebige Familien $\mathcal F,\mathcal G$ gelten
\[
\begin{aligned}
 \bigcup(\mathcal F\cap\mathcal G)
    &\subseteq(\bigcup\mathcal F)\cap(\bigcup\mathcal G),\\
 (\bigcup\mathcal F)\setminus(\bigcup\mathcal G)
    &\subseteq\bigcup(\mathcal F\setminus\mathcal G).
\end{aligned}
\]
Die Inklusionen sind im allgemeinen keine Gleichheiten.
Die erste ist auch in \cite{paulson} verzeichnet.
\end{theorem}

\section{Durchschnitt einer nichtleeren Familie}

\begin{definition}[Absoluter Familiendurchschnitt]
Für $\mathcal F\ne\varnothing$ wird definiert
\[
 \bigcap\mathcal F
   =\{x\in\bigcup\mathcal F\mid
                         \forall A\in\mathcal F\,(x\in A)\}.
\]
Dann gilt das uneingeschränkte Elementkriterium
\[
 x\in\bigcap\mathcal F
        \Longleftrightarrow\forall A\in\mathcal F\,(x\in A).
\]
Der Ausdruck $\bigcap\varnothing$ bleibt undefiniert:
Eine Menge mit diesem uneingeschränkten Elementkriterium existiert
für die leere Familie nicht. \cite{pila}
\end{definition}

\begin{theorem}[Durchschnitt und Enthaltensein]
Für $\mathcal F\ne\varnothing$ gelten
\[
\begin{gathered}
 \bigcap\mathcal F\subseteq A\quad(A\in\mathcal F),\\
 U\subseteq\bigcap\mathcal F
       \Longleftrightarrow\forall A\in\mathcal F\,(U\subseteq A),\\
 \bigcap\{A\}=A,\qquad\bigcap\{A,B\}=A\cap B.
\end{gathered}
\]
Ist $\varnothing\in\mathcal F$, so ist
$\bigcap\mathcal F=\varnothing$.
\end{theorem}

\begin{theorem}[Teilfamilien und gemeinsame Durchschnitte]
Für $\varnothing\ne\mathcal F\subseteq\mathcal G$ gilt
\[
 \bigcap\mathcal G\subseteq\bigcap\mathcal F.
\]
Für $\mathcal F,\mathcal G\ne\varnothing$ gilt
\[
 \bigcap(\mathcal F\cup\mathcal G)
      =\bigcap\mathcal F\cap\bigcap\mathcal G.
\]
Für $\mathcal F\ne\varnothing$ und jede Menge $A$ gilt
$\bigcap(\mathcal F\cup\{A\})=(\bigcap\mathcal F)\cap A$.
\end{theorem}

\begin{theorem}[Verteilen über eine Mengenfamilie]
Für jede Familie $\mathcal F$ gelten
\[
\begin{aligned}
 A\cap\bigcup\mathcal F
       &=\bigcup\{A\cap B\mid B\in\mathcal F\},\\
 (\bigcup\mathcal F)\setminus A
       &=\bigcup\{B\setminus A\mid B\in\mathcal F\}.
\end{aligned}
\]
Für $\mathcal F\ne\varnothing$ gelten außerdem
\[
\begin{aligned}
 A\cup\bigcap\mathcal F
       &=\bigcap\{A\cup B\mid B\in\mathcal F\},\\
 A\cap\bigcap\mathcal F
       &=\bigcap\{A\cap B\mid B\in\mathcal F\}.
\end{aligned}
\]
Für beliebige $A,\mathcal F$ gilt die genaue Fallunterscheidung
\[
 \begin{gathered}
 A\cup\bigcup\mathcal F
       =\bigcup\{A\cup B\mid B\in\mathcal F\}\\
 \Longleftrightarrow\quad \mathcal F\ne\varnothing\ \lor\ A=\varnothing.
 \end{gathered}
\]
\end{theorem}

\section{Relativer Durchschnitt, auch für leere Familien}

Soweit nicht anders angegeben, schließen die Aussagen leere Familien ein.

\begin{definition}[Durchschnitt innerhalb einer festen Menge]
Für beliebige Mengen $U$ und beliebige Familien $\mathcal F$ sei
\[
 \bigcap\nolimits_U\mathcal F
     =\{x\in U\mid\forall A\in\mathcal F\,(x\in A)\}.
\]
Die Glieder von $\mathcal F$ müssen nicht Teilmengen von $U$ sein.
Das Zeichen $\bigcap_U$ ist ausschließlich diese Mengenschreibweise.
\end{definition}

\begin{theorem}[Leere Familie und Bezug zum absoluten Durchschnitt]
\[
 \bigcap\nolimits_U\varnothing=U,\qquad
 \bigcap\nolimits_U\{A\}=U\cap A,\qquad
 \bigcap\nolimits_U\mathcal F\subseteq U.
\]
Für $\mathcal F\ne\varnothing$ gilt
\[
 \bigcap\nolimits_U\mathcal F=U\cap\bigcap\mathcal F.
\]
\end{theorem}

\begin{theorem}[Inklusionen und Zusammenfassen relativer Durchschnitte]
\[
\begin{gathered}
 \mathcal F\subseteq\mathcal G
    \Rightarrow
       \bigcap\nolimits_U\mathcal G\subseteq\bigcap\nolimits_U\mathcal F,\\
 U\subseteq V\Rightarrow
       \bigcap\nolimits_U\mathcal F\subseteq\bigcap\nolimits_V\mathcal F,\\
 \bigcap\nolimits_U(\mathcal F\cup\mathcal G)
       =(\bigcap\nolimits_U\mathcal F)\cap
         (\bigcap\nolimits_U\mathcal G).
\end{gathered}
\]
\end{theorem}

\begin{theorem}[Verschachtelte relative Durchschnitte]
Für beliebige Mengen $U$ und Mengen $\mathfrak F$ von Familien gilt
\[
 \bigcap\nolimits_U\bigl(\bigcup\mathfrak F\bigr)
   =\bigcap\nolimits_U
            \{\bigcap\nolimits_U\mathcal F\mid\mathcal F\in\mathfrak F\}.
\]
\end{theorem}

\begin{theorem}[Distributivgesetze für zwei beliebige Familien]
Für beliebige Familien $\mathcal F,\mathcal G$ und jede Menge $U$
gelten
\[
\begin{aligned}
 (\bigcup\mathcal F)\cap(\bigcup\mathcal G)
   &=\bigcup\{A\cap B\mid A\in\mathcal F,\ B\in\mathcal G\},\\
 (\bigcap\nolimits_U\mathcal F)\cup(\bigcap\nolimits_U\mathcal G)
   &=\bigcap\nolimits_U\{A\cup B\mid A\in\mathcal F,\ B\in\mathcal G\}.
\end{aligned}
\]
\end{theorem}

\begin{theorem}[Verteilen mit angepaßtem Grundbereich]
Für beliebige Mengen $A,U$ und jede Familie $\mathcal F$ gelten
\[
\begin{aligned}
 A\cap\bigcap\nolimits_U\mathcal F
      &=\bigcap\nolimits_{A\cap U}\mathcal F,\\
 A\cup\bigcap\nolimits_U\mathcal F
      &=\bigcap\nolimits_{A\cup U}\{A\cup B\mid B\in\mathcal F\}.
\end{aligned}
\]
\end{theorem}

\begin{theorem}[De Morgansche Gesetze für beliebige Familien]
Für jede Menge $U$ und jede Familie $\mathcal F$ gelten
\[
\begin{aligned}
 U\setminus\bigcup\mathcal F
   &=\bigcap\nolimits_U\{U\setminus A\mid A\in\mathcal F\},\\
 U\setminus\bigcap\nolimits_U\mathcal F
   &=\bigcup\{U\setminus A\mid A\in\mathcal F\}.
\end{aligned}
\]
\end{theorem}

\begin{theorem}[Disjunktheit gegenüber einer Familie]
Für jede Menge $A$ und jede Familie $\mathcal F$ gilt
\[
 A\cap\bigcup\mathcal F=\varnothing
 \Longleftrightarrow
     \forall B\in\mathcal F\,(A\cap B=\varnothing).
\]
\end{theorem}
